\documentclass[11pt]{article}

%%%%%%%%%%%%%%%%%%%%%%%%%%%%%%%%%%%%%%%%%%%%%%%%%%%%%%%%
% jzexam — LaTeX Exam Package
% Author: Jorge I. Zuluaga (C) 2026-present
% Repository: https://github.com/seap-udea/jzexam
%%%%%%%%%%%%%%%%%%%%%%%%%%%%%%%%%%%%%%%%%%%%%%%%%%%%%%%%

%%%%%%%%%%%%%%%%%%%%%%%%%%%%%%%%%%%%%%%%%%%%%%%%%%%%%%%%
% EXAM PACKAGE
%%%%%%%%%%%%%%%%%%%%%%%%%%%%%%%%%%%%%%%%%%%%%%%%%%%%%%%%
% Optional: hide answers (student version)
%\usepackage[noanswers]{jzexam}
% Show answers (answer-key / instructor version)
\usepackage[answers]{jzexam}

%%%%%%%%%%%%%%%%%%%%%%%%%%%%%%%%%%%%%%%%%%%%%%%%%%%%%%%%
% OTHER PACKAGES
%%%%%%%%%%%%%%%%%%%%%%%%%%%%%%%%%%%%%%%%%%%%%%%%%%%%%%%%
\usepackage[utf8]{inputenc}
\usepackage[english]{babel}
\usepackage{amsmath}
\setlength\parindent{0pt}

%%%%%%%%%%%%%%%%%%%%%%%%%%%%%%%%%%%%%%%%%%%%%%%%%%%%%%%%
% EXAM HEADER CONFIGURATION
%%%%%%%%%%%%%%%%%%%%%%%%%%%%%%%%%%%%%%%%%%%%%%%%%%%%%%%%
\examinfo
  {Exam Template — Course X}
  {Department of Physics — University Y\\Course: General Topic\\Semester 2026-1}

% Customizing the student fill-in fields:
% - \examfield[0.5]{...}  uses 50% of the line width (two fields per row)
% - \examfield[1]{...}    uses 100% of the line width (one field per row)
% - \examfield{...}       auto-width (shares remaining space equally)
\setexamheaderfields{%
 \examfield[0.5]{Name}%
 \examfield{ID Number}%
 \examfield{Programme}%
}

% Optional: remove the initial horizontal rule from the header
%\hideinitialhruler
\begin{document}
\makeexamheader

%%%%%%%%%%%%%%%%%%%%%%%%%%%%%%%%%%%%%%%%%%%%%%%%%%%%%%%%
% PROBLEMS
%%%%%%%%%%%%%%%%%%%%%%%%%%%%%%%%%%%%%%%%%%%%%%%%%%%%%%%%
\begin{examproblems}

%%%%%%%%%%%%%%%%%%%%%%%%%%%%%%%%%%%%%%%%%%%%%%%%%%%%%%%%
% PROBLEM
%%%%%%%%%%%%%%%%%%%%%%%%%%%%%%%%%%%%%%%%%%%%%%%%%%%%%%%%
%\problempage
\problempage[header]
\begin{problem}{Generic Problem in Relativistic Mechanics}[30]
State the general context of the problem here.

\begin{questions}

%===========================
% QUESTION
%===========================
\question Derive a symbolic expression for a physical quantity of interest.

\begin{solution}
This is a sample solution. Replace this text with the actual derivation.
\end{solution}

%===========================
% QUESTION (example of \hidequestion)
%===========================
% Use \hidequestion instead of commenting out a sub-question to temporarily
% suppress it. Neither the question text nor its solution will appear in the PDF.
\hidequestion Discuss the physical interpretation of the previous result and its limiting cases.

\begin{solution}
Include a brief interpretation and a discussion of limiting cases here.
\end{solution}

\end{questions}
\end{problem}

%%%%%%%%%%%%%%%%%%%%%%%%%%%%%%%%%%%%%%%%%%%%%%%%%%%%%%%%
% PROBLEM
%%%%%%%%%%%%%%%%%%%%%%%%%%%%%%%%%%%%%%%%%%%%%%%%%%%%%%%%
\problempage[header]
\begin{problem}{Generic Problem in Electromagnetic Fields}[40]
Consider a scenario with fields \(\vec{E}\) and \(\vec{B}\) under a Lorentz transformation.

\begin{questions}

%===========================
% QUESTION
%===========================
\question Calculate an invariant quantity associated with the field tensor.

\begin{solution}
Include the invariant expression and the main derivation steps here.
\end{solution}

%===========================
% QUESTION
%===========================
\question Determine the fields in a moving reference frame and compare with the original frame.

\begin{solution}
Write the transformation equations and a qualitative comparison of results here.
\end{solution}

\end{questions}
\end{problem}

%%%%%%%%%%%%%%%%%%%%%%%%%%%%%%%%%%%%%%%%%%%%%%%%%%%%%%%%
% PROBLEM: MULTIPLE CHOICE
%%%%%%%%%%%%%%%%%%%%%%%%%%%%%%%%%%%%%%%%%%%%%%%%%%%%%%%%
%\problempage[header]
\begin{problem}{Single-Answer and Multiple-Answer Questions}[30]
Select the correct option(s) and provide a brief justification.

\begin{questions}

%===========================
% QUESTION: SINGLE CHOICE
%===========================
\singlechoicequestion{For a free particle, the conserved quantity associated with time-translation symmetry is:}

\begin{choices}[2]
\choice{The total angular momentum.}
\correctchoice{The total energy.}
\choice{The effective electric charge.}
\choice{The local scalar curvature.}
\end{choices}

%===========================
% QUESTION: MULTIPLE CHOICE
%===========================
\multiplechoicequestion{In special relativity, which of the following statements are correct?}

\begin{choices}[2]
\correctchoice{The speed of light in vacuum is invariant.}
\choice{Proper time is the same for all inertial observers.}
\correctchoice{Events simultaneous in one frame may not be simultaneous in another.}
\choice{Rest mass depends on the reference frame.}
\end{choices}

\end{questions}
\end{problem}

%%%%%%%%%%%%%%%%%%%%%%%%%%%%%%%%%%%%%%%%%%%%%%%%%%%%%%%%
% SHORT CONCEPTUAL QUESTIONS
%%%%%%%%%%%%%%%%%%%%%%%%%%%%%%%%%%%%%%%%%%%%%%%%%%%%%%%%
%\problempage[header]
\begin{problem}{Short Conceptual Questions}[30]
Answer briefly and with justification.

\begin{questions}

%===========================
% SUBSECTION: CONCEPTS
%===========================
\questionsection{General Concepts}
Describe in two or three sentences the difference between invariant quantities and observer-dependent quantities.

\begin{solution}
An invariant keeps its value under Lorentz transformations, whereas an observer-dependent quantity changes with the reference frame.
\end{solution}

%===========================
% SUBSECTION: APPLICATION
%===========================
\questionsection{Application}
Propose an example where the choice of reference frame simplifies the analysis of a problem.

\begin{solution}
A typical example is working in the centre-of-mass frame for relativistic collisions, which simplifies the conservation of energy and momentum.
\end{solution}

\end{questions}
\end{problem}

%%%%%%%%%%%%%%%%%%%%%%%%%%%%%%%%%%%%%%%%%%%%%%%%%%%%%%%%
% END OF PROBLEMS
%%%%%%%%%%%%%%%%%%%%%%%%%%%%%%%%%%%%%%%%%%%%%%%%%%%%%%%%
\end{examproblems}

%%%%%%%%%%%%%%%%%%%%%%%%%%%%%%%%%%%%%%%%%%%%%%%%%%%%%%%%
% FORMULA SHEET (optional)
%%%%%%%%%%%%%%%%%%%%%%%%%%%%%%%%%%%%%%%%%%%%%%%%%%%%%%%%
% Use this environment to generate a formula page at the end of the exam.
\begin{formulasheet}
    \item Euler--Lagrange equations: $\frac{d}{dt}\left(\frac{\partial L}{\partial \dot{q}_i}\right) - \frac{\partial L}{\partial q_i} = 0$
    
    \item Lagrange--Jacobi identity and Virial theorem: $\dot{G}=2 K+U$, $2\langle K\rangle=-\langle U\rangle$
    
    \item Lorentz transformations (1D): $x' = \gamma (x - vt), \quad t' = \gamma \left(t - \frac{vx}{c^2}\right)$
\end{formulasheet}

%%%%%%%%%%%%%%%%%%%%%%%%%%%%%%%%%%%%%%%%%%%%%%%%%%%%%%%%
% END OF DOCUMENT
%%%%%%%%%%%%%%%%%%%%%%%%%%%%%%%%%%%%%%%%%%%%%%%%%%%%%%%%
\end{document}
